\documentclass[12pt]{article}
\usepackage[utf8]{inputenc}
\usepackage[T1]{fontenc}
\usepackage{hyperref}
\title{Kw'adza Language Resources for K-12 Curriculum Development}
\author{Compiled by Claude Sonnet 5.5 (claude-sonnet-5-5) for Andr\'as Kornai}
\date{}
\begin{document}
\maketitle

\section{Introduction}
Kw'adza (ISO 639-3 wka; Glottocode kwad1248), also spelled Qwadza and sometimes called Ngomvia, is an extinct South Cushitic language of Tanzania. The last speaker is reported to have died between 1976 and 1999. It is poorly attested, and its grammar was never recorded in detail. A school curriculum would have to be a heritage and revitalization programme built on very limited records.

\section{Primary Educational Resources}
None found. No primers, textbooks or school programmes exist in the search results.

\section{Linguistic and Reference Materials}
\begin{itemize}
\item Overview of the language and its documentation: \url{https://en.wikipedia.org/wiki/Kw%CA%BCadza_language}
\item Overview of the Kw'adza people: \url{https://en.wikipedia.org/wiki/Kw%27adza_people}
\item Glottolog entry: \url{https://glottolog.org/resource/languoid/id/kwad1248}
\item The main scholarly source reported is Christopher Ehret's work, including an unpublished ``Kw'adza vocabulary'' manuscript and his 1980 Southern Cushitic reconstruction. No URL for either was seen.
\end{itemize}

\section{Religious and Community Resources}
No scripture translation or community organisation was found. Whether descendants of the Kw'adza identify as such is reported as unknown.

\section{Digital and Audio Resources}
No audio, video or corpus resources were found.

\section{Curriculum Development Context}
Nothing found supports a language-medium curriculum. The only feasible approach is a heritage module on the history and culture of the Kw'adza and their neighbours, plus a small vocabulary drawn from Ehret's material if it can be obtained. This would need consultation with descendant communities in Tanzania. The related extinct language Asa (Aasax) has only two early vocabulary lists.

\end{document}